\chapter{Kesimpulan dan Saran}

\section{Kesimpulan}

Skripsi ini berawal dari membangun sistem rekomendasi klinis untuk penyakit
jantung koroner kronis yang bersifat FHIR-native dan berbasis \emph{knowledge graph}. Pengukuran kemudian digunakan untuk mengetahui apa yang sesungguhnya diperoleh dari rancangan FHIR-native
tersebut. Sistem berhasil dibangun dan ketiga tujuan penelitian yang dinyatakan pada Bab~I kini
dapat dijawab dari hasil-hasil Bab~IV.

\begin{enumerate}
    \item \textbf{Alurnya terwujud dan terevaluasi secara menyeluruh.} \emph{Bundle} HL7
    FHIR~R4 dari 5.490 pasien dimuat ke dalam graf properti Neo4j dan disatukan dengan
    panduan AHA/ACC CCD 2023 beserta korpus literatur yang di-\emph{seed} dari referensi
    panduan klinis tersebut melalui jembatan konsep UMLS sehingga terbentuk graf berisi 1.804.379
    simpul dan 8.796.425 relasi dengan 25 label, 51 tipe relasi, serta 26 sistem kode
    (Subbab~\ref{sec:kg-results}). Cacah simpul \texttt{:Patient}-nya sama dengan jumlah
    pasien yang dilaporkan Synthea pada folder pembangkitan sehingga pemuatan idempoten tersebut tidak membuang
    maupun menggandakan satu pun rekam medis. Terhadap standar rujukan yang dibangkitkan
    melalui jalur data yang terpisah dari sistem, keselarasan panduan klinis terhadap kohort
    evaluasi 876
    pasien mencapai $F_1 = 0{,}810$ (Subbab~\ref{sec:ablation}); bila 15 rekomendasi yang
    predikatnya dipakai bersama oleh kedua sisi disingkirkan, 53 rekomendasi yang tersisa
    mencatat $F_1 = 0{,}797$ (Subbab~\ref{sec:shared-predicate}). Angka
    tersebut bertumpu pada rujukan yang tervalidasi pada 30
    pasien yang dilabeli manual, dengan standar otomatisnya sepakat pada 99,5\%
    sel (2.150 dari 2.160) dengan Cohen's
    $\kappa = 0{,}991$ (interval kepercayaan 95\% 0,984---0,996, dengan pasien sebagai satuan
    pencuplikan ulang; Subbab~\ref{sec:kappa-ci}). Ketiga puluh pasien cukup untuk mendukung
    klaim kesepakatan hampir sempurna, tetapi pelabelnya tunggal
    (Subbab~\ref{sec:limitations}). Telaah klinisi menambahkan penilaian yang tidak dapat
    diberikan standar rujukan: seorang SpJP menilai 77,7\% dari 247 rekomendasi bagi 15 pasien
    tepat (Subbab~\ref{sec:clinician}). Di luar rekomendasi antitrombotik dan penyekat beta,
    hampir seluruh rekomendasi dinilai tepat. Penilaian yang belum tepat dan ketidaksepakatan
    antarpenilai ($\kappa = 0{,}29$) terpusat pada rekomendasi antitrombotik, yang pada
    panduannya bergantung pada waktu sejak sindrom koroner akut, risiko perdarahan, dan
    indikasi antikoagulan. Dengan demikian, tujuan~1 tercapai.

    \item \textbf{\emph{Observation gate} meningkatkan keselarasan panduan klinis dan ablasinya
    melokalkan peningkatan itu secara persis.} Penggerbangan rekomendasi bergantung-observasi
    menurut nilai terukur pasien menaikkan $F_1$ dari 0,743 menjadi
    0,810, yaitu kenaikan $+0{,}067$ yang diperoleh melalui presisi
    (0,612 menjadi 0,684) dengan \emph{recall} yang praktis
    tidak berubah (Subbab~\ref{sec:ablation}). Klasifikasi pada
    Subbab~\ref{sec:disc-class} mengajukan prediksi yang dapat difalsifikasi mengenai di mana
    perolehan itu semestinya berada dan prediksi tersebut terpenuhi dengan tepat: 53 dari 68
    rekomendasi identik pada kedua lengan hingga ke cacah \emph{true positive}, \emph{false
    positive}, dan \emph{false negative}-nya; keseluruhan selisihnya berasal dari 15
    rekomendasi yang memang bergerbang. Hasil kedua yang kurang kentara adalah bahwa
    penggerbangan bukan satu mekanisme, melainkan dua. Pada lima rekomendasi \emph{tergerbang-nilai}, penggerbangan bekerja pada presisi, yang naik
    dari 0,194 menjadi 1,000 karena rekomendasi itu tidak lagi menyala bagi setiap anggota
    kohort. Pada sepuluh rekomendasi kelas \emph{keduanya}, penggerbangan justru bekerja
    terutama pada \emph{recall}, yang naik dari 0,006 menjadi 1,000. Penyebabnya, target
    kondisi terkode milik rekomendasi itu jarang cocok bila berdiri sendiri sehingga
    predikat numeriknyalah yang mengenali pasien yang dikenai rekomendasi tersebut
    (Subbab~\ref{sec:ablation-class}). Karena itu, menceritakan ablasi ini
    semata-mata sebagai peningkatan presisi---sebagaimana terlihat oleh angka agregatnya---
    keliru bagi dua pertiga rekomendasi yang terpengaruh. Satu kualifikasi melekat pada keseluruhan ablasi ini. Kelima belas rekomendasi yang
    tergerbang itu justru merupakan rekomendasi yang predikatnya dipakai bersama oleh sistem
    dan standar rujukannya. Karena itu, ablasi ini mengukur kesetiaan implementasi, bukan
    keabsahan klinis ambangnya (Subbab~\ref{sec:shared-predicate}). Tujuan~2 tercapai.

    \item \textbf{\emph{Faithfulness} pada tingkat sitasi dapat diputuskan secara mekanis,
    tetapi \emph{faithfulness} makna tidak terjamin.}
    \emph{Grounded briefing} menelusuri setiap faktanya secara deterministik dari graf,
    menyertakan pengenal rekomendasi, kelas rekomendasi, dan PMID pada tiap fakta, serta
    membatasi model bahasa hanya pada perumusan kalimat. Seluruh 30 ringkasan yang dibangkitkan
    lolos pemeriksaan otomatis: nol rekomendasi atau PMID yang dikarang, nol rekomendasi yang
    terlewat, dan nol klaim antikoagulan yang berlebih (Subbab~\ref{sec:faithfulness}).
    Bagian yang struktural dari hasil itu adalah bahwa model berada di hilir penelusuran tanpa
    jalan kembali ke graf sehingga rekomendasi di luar hasil penelusuran pasti karangan dan
    pasti terdeteksi. Pemeriksaan butir yang dapat disitasi karenanya dapat diputuskan, sesuatu
    yang tidak dapat ditawarkan \emph{retrieval-augmented generation} berbasis kemiripan.
    Penelusuran yang deterministik tidak menjamin \emph{faithfulness} makna: model masih dapat
    merumuskan fakta yang benar secara keliru dan pemeriksaan ini hanya menangkap satu pola
    kekeliruan semacam itu. Pola itu pun diperbaiki sesudah keluarannya terlihat. Tujuan~3
    tercapai dalam batas tersebut.
\end{enumerate}

Di luar ketiga tujuan tersebut, muncul satu temuan yang tidak diantisipasi sejak awal dan
berdampak pada setiap penerapan berbasis FHIR. Dua proyeksi dari rekam medis sintetis yang
sama ternyata tidak sebanding kekayaannya. Ekspor FHIR mempertahankan status pengobatan yang
kategoris, tetapi membuang tanggal berhenti, sedangkan ekspor CSV mempertahankan tanggal
berhenti, tetapi tidak membawa status. Tidak ada satu pun di antaranya yang merupakan
subhimpunan yang lain. Aturan NSAID kronis bergantung pada berapa lama sebuah
obat dikonsumsi sehingga hanya dua dari empat cabang aturannya yang dapat dihitung dari graf
(Subbab~\ref{sec:chronicity}). Akibatnya, varian keamanan
tergerbang-kronisitas mencapai presisi 0,909 pada \emph{recall}
0,526 dan yang terlewatkan merupakan pasien yang kekronisannya hanya
terbentuk melalui satu periode penggunaan yang panjang
(Subbab~\ref{sec:safety-detection}). Sistem yang hanya mengonsumsi FHIR wajib
mengantisipasi jenis kehilangan informasi ini di setiap tempat panduan klinisnya yang bergantung pada
durasi paparan.

Asimetri tersebut berlaku pula pada arah sebaliknya dan hal itu menyempitkan satu klaim yang
semula dinyatakan penelitian ini. Perolehan presisi \emph{observation gate} diatributkan kepada
rancangan FHIR-native sejak Bab~I, tetapi atribusi itu semula dinyatakan tanpa diuji. Ketika
gerbang yang sama dibangun kembali dengan hanya membaca ekspor tabular, tidak satu pun dari
13.140 keputusan yang diperbandingkan berubah dan kohortnya pun terpilih sama persis
(Subbab~\ref{sec:csv-gating}). Perolehannya sendiri nyata dan tidak berubah. Sumbernya yang
keliru dinyatakan: yang menghasilkannya adalah pemuatan nilai observasi numerik sebagai data
kelas satu dan FHIR merupakan cara yang paling langsung untuk memperolehnya, bukan satu-satunya
cara. Kesimpulan ini diperoleh dengan menguji klaim sendiri alih-alih mempertahankannya dan
temuan negatif seperti ini justru perlu dilaporkan karena atribusi semacam itu lazim dinyatakan
tanpa pengujian pada penelitian sejenis.

\section{Saran}

Saran-saran berikut mengikuti keterbatasan yang diuraikan pada Subbab~\ref{sec:limitations},
diurutkan menurut seberapa besar masing-masing akan mengubah hasil penelitian.

\begin{enumerate}
    \item \textbf{Melebarkan gerbang kohort menjadi disjungsi konsep.} Gerbang CCD hanya
    menerima satu konsep UMLS, sedangkan panduan klinisnya mendefinisikan penyakit jantung koroner
    kronis secara disjungtif. Terdapat 67 pasien---$7{,}1\%$ dari kohort
    yang didefinisikan secara klinis---yang bukti CCD-nya semata berupa kode
    \emph{situation} infark terdahulu atau revaskularisasi terdahulu karenanya tidak pernah
    dicocokkan dengan satu pun rekomendasi dan pada penerapan nyata mereka tidak akan
    menerima keluaran apa pun (Subbab~\ref{sec:gate}). Perubahan ini kecil, hanya menyentuh
    satu kueri, dan merupakan perolehan cakupan terbesar yang tersedia.

    \item \textbf{Memperluas telaah klinisi dan menyelenggarakan studi evaluasi \emph{briefing}.}
    Telaah klinisi pada penelitian ini hanya mencakup 15 pasien bagi satu SpJP dan kesepakatan
    antarpenilainya tergolong cukup (Subbab~\ref{sec:clinician}). Telaah lanjutan sebaiknya
    melibatkan beberapa SpJP yang menilai seluruh pasien. Hasil telaah juga menunjuk perbaikan
    pada sistemnya. Kesimpulan \emph{briefing} perlu menyebutkan obat yang ditandai serta
    informasi yang menjadi dasar rekomendasi antitrombotik dan penyekat beta, misalnya waktu
    sejak sindrom koroner akut, fraksi ejeksi, dan status gagal jantung. Bila informasi itu
    tidak dapat dipastikan dari rekam medis, seperti risiko perdarahan, kesimpulan perlu
    menyatakannya. Ketiadaan informasi itulah yang paling sering menghalangi penilaian. Selain itu,
    \emph{faithfulness} merupakan sifat hubungan ringkasan dengan rekam medis, bukan sifat
    kebergunaannya bagi pembaca. Studi dengan pola Van~Veen \textit{et al.}~\cite{vanveen2023},
    yang di dalamnya klinisi menilai ringkasan atas kelengkapan, kebenaran, dan keringkasannya
    dengan keandalan antarpenilai diukur melalui Cohen's $\kappa$ \cite{cohen1960}, akan
    menutup kesenjangan tersebut. Studi itu sekaligus mengangkat evaluasi \emph{faithfulness}
    dari pemeriksaan otomatis atas butir yang dapat disitasi menjadi penilaian
    \emph{faithfulness} makna oleh pembaca manusia.

    \item \textbf{Mengulang evaluasi atas rekam medis nyata.} Pasien Synthea dibangkitkan
    dari model perkembangan penyakit sehingga pola komorbiditas dan lintasan
    laboratoriumnya konsisten secara internal, tetapi tidak diambil dari populasi nyata. Hasil
    yang dilaporkan di sini memastikan bahwa alurnya benar secara internal, tetapi hasil itu belum
    memastikan bahwa capaian keselarasan yang sama akan diperoleh pada rekam medis yang lebih
    berderau, lebih jarang, dan disertai teks tak terstruktur. Rekam medis nyata juga membawa
    gejala dan keluhan jantung yang tidak dimuat Synthea sehingga enam rekomendasi yang
    bergantung pada angina baru dapat diuji di sana (Subbab~\ref{sec:limitations}).

    \item \textbf{Memulihkan logika panduan klinis yang bergantung pada durasi.} Profil FHIR yang
    membawa periode penyerahan atau periode dosis, ataupun pemasangan graf dengan sumber
    tambahan, akan membuat keempat cabang aturan penggunaan kronis dapat dihitung sekaligus
    menghapus batas atas \emph{recall} yang terukur pada
    Subbab~\ref{sec:safety-detection}.

    \item \textbf{Memberi lapisan keamanan dengan representasi yang lebih kaya.} Pencocokan tergerbang-populasi memiliki presisi sempurna, tetapi hanya memulihkan sebagian
    kecil kontraindikasi, yaitu yang kondisi populasinya dapat digerbangi secara struktural.
    Kontraindikasi yang terlewat adalah yang didefinisikan lewat negasi dan yang bersifat
    konjungtif. Keduanya tidak dapat diungkapkan oleh skema datar
    (Subbab~\ref{sec:safety-detection}). Merepresentasikan keduanya secara eksplisit,
    alih-alih memilih antara penandaan berlebih dan pemulihan yang kurang, merupakan jalan
    keluar dari pertukaran tersebut.

    \item \textbf{Memindahkan rancangan ini ke panduan klinis atau penyakit yang lain.} Hanya
    panduan AHA/ACC CCD 2023 yang dikodekan dan hanya untuk penyakit jantung koroner
    kronis. Rancangan jembatan konsepnya membuat generalisasi masuk akal karena tidak ada
    bagian lapisan kueri yang khas kardiologi, tetapi penelitian ini belum membuktikan kemampuan tersebut.
    Mengodekan panduan klinis kedua lalu menjalankan ulang evaluasi yang sama akan
    membuktikannya.
\end{enumerate}
